\documentclass[11pt,a4paper]{article}
\usepackage[margin=24mm,headheight=15pt]{geometry}
\usepackage{fontspec}
\setmainfont{Times New Roman}
\setsansfont{Avenir Next}
\setmonofont{Menlo}
\usepackage{microtype}
\usepackage{newtxmath}
\usepackage{mathtools}
\usepackage{booktabs,longtable,array,tabularx}
\usepackage{pdflscape}
\usepackage[table]{xcolor}
\usepackage{enumitem}
\usepackage{fancyhdr}
\usepackage{titlesec}
\usepackage{xurl}
\usepackage{hyperref}
\usepackage{bookmark}
\usepackage{fancyvrb}
\usepackage[most]{tcolorbox}

\definecolor{Navy}{HTML}{18324A}
\definecolor{Accent}{HTML}{C56A3A}
\definecolor{LinkBlue}{HTML}{2D6A8A}
\definecolor{SoftBlue}{HTML}{EAF1F5}
\definecolor{TableHead}{HTML}{DCE8EF}
\definecolor{BodyGray}{HTML}{26343D}

\hypersetup{
  colorlinks=true,
  linkcolor=LinkBlue,
  urlcolor=LinkBlue,
  citecolor=LinkBlue,
  pdftitle={Comparison and research recommendation},
  pdfauthor={Final Thesis Master Project}
}
\setlength{\parindent}{0pt}
\setlength{\parskip}{0.65em}
\setlength{\emergencystretch}{3em}
\setlist{leftmargin=1.6em,itemsep=0.3em,topsep=0.25em}
\setcounter{tocdepth}{2}
\renewcommand{\arraystretch}{1.15}
\allowdisplaybreaks

\titleformat{\section}
  {\Large\bfseries\sffamily\color{Navy}}
  {\thesection}{0.65em}{}
  [\vspace{-0.25em}\color{Accent}\titlerule]
\titleformat{\subsection}
  {\large\bfseries\sffamily\color{Navy}}
  {\thesubsection}{0.6em}{}
\titleformat{\subsubsection}
  {\normalsize\bfseries\sffamily\color{LinkBlue}}
  {\thesubsubsection}{0.55em}{}

\newtcolorbox{ResearchQuote}{
  enhanced,
  breakable,
  colback=SoftBlue,
  colframe=LinkBlue,
  boxrule=0pt,
  leftrule=2.2pt,
  arc=0pt,
  left=8pt,right=8pt,top=5pt,bottom=5pt
}
\newenvironment{ResearchContents}
  {\begin{tcolorbox}[breakable,colback=SoftBlue,colframe=SoftBlue,arc=2pt]}
  {\end{tcolorbox}}

\pagestyle{fancy}
\fancyhf{}
\fancyhead[L]{\small\sffamily\color{Navy} Asymmetric stochastic TSP}
\fancyhead[R]{\small\sffamily\color{Navy}\detokenize{comparison.md}}
\fancyfoot[C]{\small\sffamily\color{BodyGray}\thepage}
\renewcommand{\headrulewidth}{0.3pt}
\renewcommand{\footrulewidth}{0pt}

\begin{document}
\color{BodyGray}
\begin{titlepage}
  \thispagestyle{empty}
  \vspace*{22mm}
  {\sffamily\small\bfseries\color{Accent}\MakeUppercase{Research note}}\par
  \vspace{8mm}
  {\sffamily\fontsize{21}{25}\selectfont\bfseries\color{Navy} Comparison and research recommendation}\par
  \vspace{6mm}
  {\color{Accent}\rule{42mm}{1.4pt}}\par
  \vspace{11mm}
  {\Large\color{BodyGray} Asymmetric stochastic TSP\\
  clairvoyance-gap investigation}\par
  \vfill
  \begin{tcolorbox}[
    colback=SoftBlue,colframe=SoftBlue,arc=2pt,
    left=10pt,right=10pt,top=8pt,bottom=8pt]
    \sffamily
    \textbf{Source}\quad \detokenize{comparison.md}\\[3pt]
    \textbf{Prepared}\quad 25 July 2026\\[3pt]
    \textbf{Format}\quad Typeset mathematical PDF
  \end{tcolorbox}
  \vspace{12mm}
\end{titlepage}

\begin{ResearchContents}
\tableofcontents
\end{ResearchContents}
\clearpage

\phantomsection
\section*{Bottom line}
\addcontentsline{toc}{section}{Bottom line}

A hybrid of the switching-network and design viewpoints gives a complete construction above \(4/3\).  The four-star layered poset in \nolinkurl{four-star-layered-poset-construction.md} has


\[
 \mathbb E W_{\rm post}=2+o(1),\qquad
 \inf_\pi\mathbb E K_\pi\ge3-o(1),
\]

and its positive directed-metric embedding proves


\[
 \liminf
 \frac{\operatorname{OPT}_{\rm adapt}}
      {\operatorname{OPT}_{\rm post}}\ge\frac32.
\]

The finite certificate in \nolinkurl{four-star-poset-finite-certificate.md} gives one explicit strict instance without computation.


The key successful feature is not expansion or girth.  Two permanent gate lanes force every two-run causal execution to make a local \(K_{2,2}\) edge commitment.  Four independent row/column-star selectors make one exact two-active pattern fatal for every local decision rule.  A Bellman product supermartingale multiplies this loss across layers even under arbitrary remote inactive probes.


\phantomsection
\section*{Status table}
\addcontentsline{toc}{section}{Status table}

\begin{landscape}

\begingroup\small
\setlength{\tabcolsep}{3pt}
\renewcommand{\arraystretch}{1.25}
\begin{longtable}{>{\raggedright\arraybackslash}p{0.15\linewidth}>{\raggedright\arraybackslash}p{0.27\linewidth}>{\raggedright\arraybackslash}p{0.28\linewidth}>{\raggedright\arraybackslash}p{0.14\linewidth}}
\rowcolor{TableHead}
\textbf{Family} & \textbf{Strongest justified result} & \textbf{Main obstruction} & \textbf{Priority} \\
\hline
\endfirsthead
\rowcolor{TableHead}
\textbf{Family} & \textbf{Strongest justified result} & \textbf{Main obstruction} & \textbf{Priority} \\
\hline
\endhead
Layered four-star switching/design hybrid & \textbf{Complete theorem:} positive asymmetric metric with gap liminf at least \(3/2\), plus an explicit finite \(>4/3\) witness. & No obstruction in the proved construction; the Bellman product is needed to handle interleaving. & \textbf{Solved positive direction.} \\
\hline
Projective-plane incidence graphs & Exact theorem: the uniform point/line directed template on \(\mathrm{PG}(2,q)\) has gap \(1\). & The asymmetry is a vertex potential; a Singer Hamilton cycle gives an activation-neutral order. & Stop for the uniform template. \\
\hline
Generalized polygons & Exact potential no-go; exact gap \(1\) on balanced Hamiltonian subclasses. & Uniform incidence direction is potential; finite thick apartments have bounded length by Feit--Higman. & Low. Use only if a nonpotential circulation is explicit. \\
\hline
Cayley digraphs & Exact finite Cayley-metric gadget with gap \(8/7\). & Noncommutativity separates complete words but does not defeat selector-first calls or translated local repairs. & Medium. Best candidate host for a proved route cell. \\
\hline
High-girth lifts & Exact articulation-lift additivity: a high-girth scaffold carrying independent Chapter 4 modules stays below \(4/3\). & Girth constrains local paths, not remote calls; port coupling also creates posterior circulation defects. & Medium-low until a base route cell exists. \\
\hline
Ramanujan bigraphs & Proved realization-wise Euler-sweep upper bound and a legal causal pending sweep with the same bound; exact divergence calculation for the all-preferred-traces ledger. & Shared endpoints destroy additive charges; preferred orientations are imbalanced; inactive edges need not be used. & Medium-low. Only the sparse \(p=\lambda/d\) star-routing version remains credible. \\
\hline
Linear codes, designs, LTCs & Exact legality/no-go results plus a conditional regular-design theorem with a sharp \(4/3\) threshold. & Code distance is combinatorial; the missing commitment and wrong-mode charges are directed-geometric. & \textbf{Highest.} This is the cleanest lower-bound framework. \\
\hline
Finite buildings & Exact potential no-go for radial galleries; apartment-size, selector, and retraction diagnostics. & Independent active chambers do not select one apartment; retractions collapse service multiplicity. & Medium-low. Affine quotient cocycles are the only distinctive opening. \\
\hline
Switching networks & Exact gap \(1\) for the natural butterfly/Beneš output-terminal template. & A static reusable fabric has no persistent hidden switch state; every metric traversal can choose a fresh route. & \textbf{High as a gadget question}, low as an off-the-shelf construction. \\
\hline
Algebraic lifts & Exact pointwise-permutation obstruction and exact random defect count. & Independent pointwise choices do not form a cycle cover; the \(7/4\) and \(8/7\) calculations are only a conditional retained-trace ledger, not optimum bounds. & Low for direct pointwise lifts; medium as a certificate inside another gadget. \\
\hline
\end{longtable}
\endgroup

\end{landscape}

\phantomsection
\section*{Strongest complete positive calculations}
\addcontentsline{toc}{section}{Strongest complete positive calculations}

The main construction proves


\[
 \liminf
 \frac{\operatorname{OPT}_{\rm adapt}}
      {\operatorname{OPT}_{\rm post}}\ge\frac32.
\]

Its smallest recorded finite certificate has \(120002\) clients and proves


\[
 3\operatorname{OPT}_{\rm adapt}
 -4\operatorname{OPT}_{\rm post}>0.24.
\]

An earlier, smaller positive calculation remains useful as a local control:


\[
 \operatorname{OPT}_{\rm post}=\frac72,\qquad
 \operatorname{OPT}_{\rm adapt}=4,\qquad
 \frac{\operatorname{OPT}_{\rm adapt}}
      {\operatorname{OPT}_{\rm post}}=\frac87
\]

for the finite Cayley word-order gadget.


The other exact calculations are no-go examples:


\begin{itemize}

\item projective-plane uniform incidence: ratio \(1\);

\item butterfly output terminals: ratio \(1\);

\item balanced Hamiltonian generalized-polygon uniform incidence: ratio \(1\).

\end{itemize}

These are useful controls: vertex transitivity does not force ratio \(1\) (the Cayley gadget disproves that), whereas high path diversity does not force a gap (the butterfly theorem disproves that).


\phantomsection
\section*{The sharp conditional route-code theorem}
\addcontentsline{toc}{section}{The sharp conditional route-code theorem}

The code/design note gives the most useful quantitative target.  Let \(k\) independent fair selector bits determine \(m\) parity route modes through a regular design.  Suppose:


\begin{itemize}

\item every realization has a posterior tour of cost at most \(B\);

\item querying a selector before route commitment costs \(a\) in expectation;

\item a wrong route mode costs \(\Delta\);

\item the baseline, early-query, and wrong-mode charges are disjoint after arbitrary metric expansion.

\end{itemize}

Then


\[
 \operatorname{OPT}_{\rm adapt}
 \ge B+\min\{ak,\Delta m/2\},
\qquad
 \frac{\operatorname{OPT}_{\rm adapt}}
      {\operatorname{OPT}_{\rm post}}
 \ge
 1+\frac{\min\{ak,\Delta m/2\}}{B}.
\]

Thus the construction beats \(4/3\) exactly when


\[
 \min\{ak,\Delta m/2\}>B/3.
\]

The single commitment cut in this theorem is deliberately strong.  An actual construction will probably allow progressive/interleaved commitments.  The right replacement is an LTC-style Bellman subsolution that tracks unresolved test risk after every possible next call.


\phantomsection
\section*{Cross-family eliminators}
\addcontentsline{toc}{section}{Cross-family eliminators}

These checks should be applied before developing another construction.


\begin{enumerate}

\item \textbf{Potential test.}  If \(w(u,v)=s_{\{u,v\}}+\phi(v)-\phi(u)\), the directed part telescopes on every closed realized call sequence.  It creates no asymmetric effect.

\item \textbf{Product-law test.}  Random codewords, lift signs, apartment choices, switch settings, and edge orientations are not hidden data.  Only independent client activations are.

\item \textbf{Inactive-transit test.}  Inactive clients and their incident arcs remain usable by shortest paths.  Activation is a service requirement, not edge availability.

\item \textbf{Selector-first test.}  The policy may remotely call every selector. The active subsequence must already cost the scale claimed in the lower bound.

\item \textbf{Additivity test.}  If post and adaptive costs decompose over modules, an outer graph cannot improve the largest local ratio.

\item \textbf{Flow-consistency test.}  Locally preferred directed traces need not form a circulation or tour.  Holes, collisions, and vertex divergence belong in the posterior accounting.

\item \textbf{Interleaving test.}  A block, apartment, fiber, line, or switch cell cannot be assumed served in one visit.  Every extra service piece needs a distinct or bounded-congestion charge.

\item \textbf{Metric-closure test.}  Expansion and high connectivity usually add repair paths.  Every claimed penalty needs a potential, quotient norm, scale separation, or port lemma.

\end{enumerate}

\phantomsection
\section*{What made the successful construction work}
\addcontentsline{toc}{section}{What made the successful construction work}

\phantomsection
\subsection*{Local delayed route cell}
\addcontentsline{toc}{subsection}{Local delayed route cell}

The two permanent gates at every layer prevent an active future probe from being harmless: jumping over an unserved gate leaves two incomparable gates for the last run and certifies a third run.  This supplies the delayed commitment that the natural butterfly and reusable-switch templates lacked.


\phantomsection
\subsection*{Second-order versus third-order scaling}
\addcontentsline{toc}{subsection}{Second-order versus third-order scaling}

The causal conflict occurs on an exact two-active event of probability \(\Theta(q^2)\), while posterior width exceeds two only on three-active events of probability \(\Theta(q^3)\).  Choosing \(Lq^2\to\infty\) and \(Lq^3\to0\) separates the two objectives.


\phantomsection
\subsection*{Interleaving-safe repetition}
\addcontentsline{toc}{subsection}{Interleaving-safe repetition}

The product proof uses the exact local Bellman success value after ordinary queries, factor \(1-h\) for an untouched future stage, and the transition


\[
 q\cdot0+(1-q)\cdot1\le1-h
\]

for its first premature query.  The product of these factors is a supermartingale, so cross-layer interleaving and free negative information do not invalidate repetition.


\phantomsection
\section*{Overall recommendation}
\addcontentsline{toc}{section}{Overall recommendation}

Use the four-star layered poset as the new baseline construction.  The most valuable next questions are whether its \(3/2\) lower bound can be sharpened, whether the client count in the finite witness can be reduced, and whether a larger \(k\times k\) star interface produces more than three forced causal runs while keeping posterior width near \(k\).

\end{document}
